\documentclass[10pt,a4paper]{article}
\usepackage[utf8]{inputenc}
\usepackage[T2A,T1]{fontenc}
\usepackage{lmodern}
\usepackage[russian,english]{babel}
\usepackage[margin=1in]{geometry}
\usepackage{amsmath,amssymb}
\usepackage{booktabs}
\usepackage{hyperref}
\usepackage{fancyhdr}
\usepackage{tabularx}

\hypersetup{
    colorlinks=true,
    linkcolor=blue,
    citecolor=blue,
    urlcolor=blue,
    pdftitle={Universal AI Filter (UAIF): A Deterministic Input-Centric Shield for Mitigating Tokenization-Level LLM Vulnerabilities},
    pdfauthor={Aleksandr Brailovsky}
}

\pagestyle{fancy}
\fancyhf{}
\rhead{\small UAIF Specification V1.0}
\lhead{\small A. Brailovsky | Studio "Architects"}
\cfoot{\thepage}

\title{\textbf{Universal AI Filter (UAIF): A Deterministic Input-Centric Shield for Mitigating Tokenization-Level LLM Vulnerabilities}}
\author{\textbf{Aleksandr Brailovsky}\\Studio "Architects" (UAIF Labs)\\
\small Contact: \href{mailto:architectsstudio54@gmail.com}{architectsstudio54@gmail.com} $|$ Repository: \href{https://github.com/uaif-specification}{github.com/uaif-specification}}
\date{July 2026 $|$ Document ID: UAIF-SPEC-2026-V1.0}

\begin{document}

\maketitle

\begin{abstract}
The reliance on probabilistic neural classifiers (Guardrails) to secure Large Language Models (LLMs) represents a fundamental architectural flaw. Attackers have shifted their focus from semantic exploits to sub-symbolic tokenization-level attacks, including Emoji Smuggling, Invisible Glyph Injection, and Homoglyph exploits, against which existing ML-based systems demonstrate a 100\% blindness rate. This paper introduces the Universal AI Filter (UAIF), a paradigm shift toward deterministic, Input-Centric Security. Operating as an isolated infrastructure proxy, UAIF enforces strict byte-level normalization, multi-layered obfuscation decoding, and visual-semantic glyph collapse prior to tokenization. This approach mathematically eliminates the attack vector by ensuring the idempotency of the input stream. Furthermore, UAIF reduces AI security operational expenditure (OPEX) by up to 66\% by migrating the security workload from expensive GPU clusters to lightweight, high-throughput CPU nodes. This manuscript serves as a foundational disclosure and prior art for open deterministic AI security standards.
\end{abstract}

\noindent\textbf{Keywords:} LLM Security, Input Normalization, Tokenization Exploits, Prior Art, Deterministic Guardrails, UAIF, Prompt Injection Defense.

\vspace{0.3cm}
\noindent\rule{\textwidth}{0.4pt}
\begin{center}
\small\textbf{License \& Defensive Publication Notice}\\
\small This specification, manuscript, and associated mathematical paradigms are released under the \textbf{MIT License} for open global dissemination, standardization, and defensive publication as Prior Art. All rights to the UAIF\textsuperscript{TM} trademark belong to Aleksandr Brailovsky / Studio "Architects".
\end{center}
\noindent\rule{\textwidth}{0.4pt}

\section{Introduction: The Crisis of Probabilistic Security}
The era of isolated chatbots, operating within secure sandboxes, has ended. By early 2026, enterprise landscapes have undergone a profound transformation toward total integration of autonomous multi-agent systems and scalable Retrieval-Augmented Generation (RAG) architectures. AI assistants now possess direct programmatic access to transactional databases, corporate document systems, email, and confidential source code repositories.

In this high-stakes environment, the compromise of a generative AI system is no longer limited to reputational damage from unethical content generation. A successful attack today means direct, uncontrolled access to enterprise trade secrets and complete takeover of critical business logic.

Despite the exponential rise in threats, current industry approaches to securing generative AI rest on a fundamental architectural error. The cybersecurity industry, driven by hardware monopolists, continues to attempt to protect a stochastic mathematical model (the base LLM) with another probabilistic neural classifier. These semantic filters, or ML-Guardrails---such as NVIDIA NeMo Guardrails, Meta Llama Guard, and Azure Prompt Shield---represent an evolutionary dead end. Entrusting one trainable neural network with the task of semantically guessing "malicious intent" in an input stream addressed to another neural network creates an endless race of probabilities that enterprise security inevitably loses.

The key problem is that modern language models conceptually inherit vulnerabilities analogous to the von Neumann architecture flaw, where executable code and processed data share the same address space. In the LLM context, system prompts (developer instructions) and unstructured, potentially hostile user input are concatenated into a single, inseparable token stream. The model, operating on probabilistic distributions to predict the next token, lacks any built-in deterministic mechanism to authenticate the source of instructions and cannot mathematically guarantee a boundary between trusted corporate patterns and attacker payloads.

Attackers and advanced persistent threat (APT) groups have recognized this weakness and long ago shifted their adversarial focus to a deeper, pre-semantic level. They no longer attempt social engineering to coax secrets from the model in natural language; they exploit the very physics of data formation---the raw byte structure and the mechanics of Byte Pair Encoding (BPE) tokenizers. At this sub-symbolic, mathematical level, semantic filters trained to detect "toxicity" or direct imperative commands are rendered completely blind.

\section{Anatomy of Sub-Symbolic Attacks: The Perception Gap}
Modern LLM architectures process input not as flat human-readable text, but via compression and tokenization algorithms that translate character sequences into discrete numerical identifiers. Attack vectors exploiting this mechanism are based on a fundamental perception gap between semantic Guardrails and the internal tokenization pipeline of the target model. Security filters analyze text as a human sees it, while the base LLM responds to its raw mathematical representation.

\subsection{Emoji Smuggling and Zero-Width Injections}
At the core of physical-layer vulnerabilities lies sophisticated manipulation of Unicode specifications. One of the most destructive vectors, which completely discredits probabilistic filters, is the Emoji Smuggling attack and the injection of invisible zero-width control characters.

Historically, Unicode variation selectors (range U+FE00--U+FE0F) and specialized tags (U+E0000--U+E007F) are designed to specify the visual style of a preceding glyph, for example, rendering a standard text character as a colorful emoji. Because these bytes have no graphical representation, they are fully transparent to visual inspection and are ignored by classifiers during string filtering.

An attacker converts a malicious system instruction into binary format, mapping each bit to a specific variation selector, and injects this array of invisible bytes immediately after a visually neutral symbol, such as a thumbs-up emoji. To a human or a security model (e.g., Azure Prompt Shield or ProtectAI), the incoming stream appears as a benign greeting with an emoji. However, the BPE tokenizer of the base LLM, lacking these characters in its compression dictionary, reads the hidden bytes and obediently expands them into a sequence of executable system directives directly inside the context window.

Research by Hackett et al. (2025) demonstrates the catastrophic failure of traditional defense systems against such methods. When using variation selectors and Unicode injections, the attack success rate reaches absolute values, completely negating multi-million-dollar corporate security investments.

\begin{table}[h!]
\centering
\caption{Guardrail System Vulnerability to Sub-Symbolic Unicode Injections}
\vspace{2mm}
\begin{tabularx}{\textwidth}{X X c c}
\toprule
\textbf{Target Guardrail System} & \textbf{Attack Vector} & \textbf{Baseline Block Rate (\%)} & \textbf{Achieved ASR (\%)} \\
\midrule
Microsoft Azure Prompt Shield & Character Injection & 28.02 & 71.98 \\
Meta Prompt Guard & Character Injection & 29.56 & 70.44 \\
NVIDIA NeMo Guard Jailbreak Detect & Character Injection & 27.46 & 72.54 \\
Azure / Meta / NVIDIA (All) & Emoji Smuggling & 0.00 & 100.00 \\
ProtectAI Prompt Injection v1/v2 & Unicode Tags Smuggling & $\sim$10.00 & 90.00 \\
\bottomrule
\end{tabularx}
\end{table}

Semantic protection, operating in terms of natural language and trained to detect toxicity or known jailbreak patterns, is mathematically incapable of recognizing a threat encrypted in the binary representation of invisible markers.

\subsection{Homoglyph Attacks and Token Fragmentation}
Homoglyph attacks involve substituting characters from one alphabet with visually identical characters from another (e.g., replacing Latin "a" with Cyrillic "\foreignlanguage{russian}{а}"). To the human eye and semantic filter, the text remains unchanged, but this completely disrupts vector database search logic and confuses the BPE tokenizer, forcing the model to perceive the substituted word as a completely new, unknown identifier absent from blacklists.

This fragmentation (TokenBreak) artificially inflates token count, leading to billing attacks (Denial of Spend) that can increase enterprise API costs by 5.2x.

\subsection{The EchoLeak Zero-Click Incident (CVE-2025-32711)}
The failure of semantic filters reaches a critical level of unacceptable risk in multi-agent RAG systems. The EchoLeak vulnerability (CVSS 9.3) in Microsoft 365 Copilot became the first documented real-world zero-click indirect prompt injection exploit. The attack required no user interaction: a malicious email with hidden instructions was automatically indexed by the RAG agent, extracted into the LLM context, and triggered silent data exfiltration via rendered Markdown images.

EchoLeak proved that protecting only user input (Input-Centric Security) is catastrophically insufficient in enterprise environments where AI agents autonomously consume, aggregate, and interpret massive volumes of unstructured data from external sources.

\section{The UAIF Paradigm: Deterministic Input-Centric Security}
Project UAIF (Universal AI Filter) implements a radical architectural shift. The core concept is a complete abandonment of attempts to interpret the semantic "meaning" or "intent" of a threat. The system operates at the preprocessing stage as a tightly isolated infrastructure proxy (the "Black Box" principle), hardware-wise and mathematically separating validation logic from the vulnerable stochastic latent space of the neural network.

The gateway relies on three proprietary concepts, forming an impenetrable layered barrier before the BPE tokenizer.

\subsection{Concept 1: Hardware-Accelerated Low-Level Normalization}
Classic Web Application Firewalls (WAF) and ML-Guardrails rely on resource-intensive string checks, embedding computation, or primitive regular expressions, all easily bypassed by modern byte-level injection methods. UAIF completely changes the paradigm: it intercepts the incoming byte stream at the network kernel level, before data is interpreted as text.

The core element is a hardware-accelerated low-level normalization module. It performs deep, surgical cleaning of encapsulated control instructions, hidden neural triggers, and invisible graphical modifiers used for LLM steganography. Any hidden attacker directives are forced into harmless informational noise, preserving legitimate business context while physically eliminating the Emoji Smuggling attack vector.

\subsection{Concept 2: Intelligent Entropy Analyzer with Circuit Breakers}
Attackers frequently use recursive multi-layer encoding (e.g., wrapping Base64 in Hex, then URL-encode, then ROT13) to hide payloads. Traditional parsers attempting recursive expansion without resource controls are vulnerable to Denial of Service (DoS) attacks---a "Zip Bomb" analog.

UAIF integrates an intelligent entropy analyzer that autonomously detects and deconstructs obfuscation of any depth under strict controls: \texttt{MaxDecodeDepth} (hard limit on recursion depth) and adaptive Circuit Breakers (automatic interruption on timeout or memory anomaly). This ensures 100\% uptime against algorithmic complexity attacks.

\subsection{Concept 3: Visual-Semantic Collapse (Mixed-Script Handling)}
Homoglyph attacks are neutralized by a proprietary Visual-Semantic Collapse mechanism. The mathematical apparatus of UAIF resolves computational collisions and blind spots in standard Unicode canonicalization libraries.

Let $I$ be an arbitrary input string subject to sub-symbolic adversarial mutations, and $\mathcal{N}$ be the UAIF normalization function executing at the byte layer:
\begin{equation}
\mathcal{N}(I) = \mathcal{N}(\mathcal{N}(I))
\end{equation}
The normalization function is fully deterministic and idempotent. It runs at the infrastructure level (CPU), enforcing a canonical representation before the input reaches the LLM tokenizer.

Importantly, UAIF incorporates an Adaptive Mixed-script analysis that dynamically recognizes legitimate language domain boundaries and applies contextual exceptions for trusted locales. This preserves the integrity of mixed Russian-English content, achieving a false positive rate (FPR) below 0.1\% on real data, while maintaining absolute resistance to targeted homoglyph exploits.

\sectionsection{Economic Value for Enterprise}
Deploying UAIF yields quantifiable benefits:
\begin{itemize}
    \item \textbf{66\% OPEX Reduction:} Security workload migrates from expensive GPU nodes to standard CPU containers.
    \item \textbf{$<$5 ms Latency Overhead:} Compatible with real-time streaming applications and high-frequency trading.
    \item \textbf{2.5 Months ROI:} Payback period based on prevented billing attacks and GPU savings.
    \item \textbf{Zero-Client Data Exposure:} On-premise, air-gapped deployment ensures full compliance with 152-FZ (Russia), ISO 27001, and ISO/IEC 42001.
\end{itemize}

\section{Conclusion}
The era of protecting artificial intelligence through intuitive neural guessing of malicious intent has reached its historical limit. Attempts to train probabilistic classifiers to recognize every variation of machine obfuscation, homoglyph substitution, and hidden Unicode injections are mathematically doomed due to infinite combinatorics and fundamental attention mechanism weaknesses.

UAIF redefines AI security by returning to the immutable principles of mathematical determinism. By integrating bidirectional RAG analysis, hardware-accelerated low-level normalization, and preventive DoS protection, UAIF denies attackers the physical space to exploit tokenization vulnerabilities. The open publication of this specification serves as a foundational prior art to accelerate the global adoption of deterministic AI security standards.

\begin{thebibliography}{9}
\bibitem{hackett2025} Hackett et al. "Bypassing LLM Guardrails: An Empirical Analysis of Evasion Attacks." \emph{arXiv:2504.11168}, 2025.
\bibitem{gao2025} Gao et al. "Imperceptible Jailbreaking against Large Language Models." \emph{arXiv:2510.05025}, 2025.
\bibitem{echoleak2025} "EchoLeak (CVE-2025-32711) -- Zero-Click Prompt Injection in Microsoft 365 Copilot." \emph{Aim Security}, 2025.
\bibitem{unicode2026} "Reverse CAPTCHA: Evaluating LLM Susceptibility to Invisible Unicode Instruction Injection." \emph{arXiv:2603.00164}, 2026.
\bibitem{firewall2025} "A Bidirectional LLM Firewall: Architecture, Failure Modes, and Evaluation." \emph{HuggingFace Discussion}, 2025.
\bibitem{canonical2025} US12278836B1 -- Canonicalization of Unicode Prompt Injections. \emph{HiddenLayer, Inc.}, 2025.
\bibitem{nist2023} NIST AI Risk Management Framework (AI RMF 1.0). \emph{National Institute of Standards and Technology}, 2023.
\end{thebibliography}

\end{document}
